\documentclass[../main.tex]{subfiles}
\begin{document}
%==========================================TIÊU ĐỀ TẬP========================================
\begin{table}[h]
  \begin{adjustwidth}{-1cm}{}
    \begin{tabular}{>{\centering\arraybackslash}p{6cm}|>{\raggedright\arraybackslash}p{12.5cm}}
      \multirow{5}{6cm}
      % Cột bên trái: ÉPISODE và số tập
      {\\[-40pt]
        \flaregothic\fontsize{20pt}{20pt}\selectfont \centering
        {\contour{errorfour}{\textcolor{black}{ÉPISODE}}}\color{black}\\
        \burbank\fontsize{72pt}{72pt}\selectfont
        \includegraphics[height=3.12359cm]{../../../../Image/Book?/number/num17.png}
        \\[18pt]
        % \flaregothic\fontsize{14pt}{14pt}\selectfont
        % \color{epattr}BIENVENUE, \uline{\color{Banpire} Mel} !
        % \flaregothic\fontsize{18pt}{18pt}\selectfont
        % \color{epattr}ÉPISODE PERDU
        \color{black}
      }
       &
      % Dòng thứ nhất: tên tập tiếng Nhật
      {\fotskip\fontsize{18pt}{18pt}\selectfont \ruby{壊}{こわ}れた\ruby{理}{り}\ruby{想}{そう}}
      \\[6pt]
      % Dòng thứ hai: tên tập tiếng Anh
       & {\cafeteria\fontsize{18pt}{18pt}\selectfont Shattered Fantasies}                         \\[4pt]
      % Dòng thứ ba: tên tập tiếng Việt
       & {\vnmsans\fontsize{18pt}{18pt}\selectfont Chấp nhận thực tại}                            \\[8pt]
      % Dòng thứ tư: Nhân vật xuất hiện
       & {\caviar{\fontsize{14pt}{14pt}\selectfont Track used: The Faded Original - Alan Walker}}
      \\[6pt]
      % Dòng cuối cùng: Thời gian diễn ra của tập (thường sẽ là ngày phát hành)
       & {\continuum\fontsize{16pt}{16pt}\selectfont Ngày phát hành: 11/06/2022}                  \\[8pt]
    \end{tabular}
  \end{adjustwidth}
\end{table}
%===========================================NỘI DUNG==========================================
\color{black}

Gió thổi qua hành lang ga cũ, mang theo mùi kim loại gỉ và bụi đất lâu năm. Ánh sáng mờ nhạt xuyên qua khung cửa kính vỡ, chiếu lên nền gạch nứt nẻ của Ga Aogami. Khi cánh cửa gỗ cũ kỹ mở ra, bốn người họ - Kuon, Kaigou, Suzuki và Shino - bước ra ngoài, lặng lẽ đối diện với thế giới vừa tái sinh sau vòng lặp.

Bên ngoài, bầu không khí khác hẳn.

Không còn tiếng tàu, không còn loa thông báo, chỉ có sự im lặng dày đặc phủ lên toàn bộ thị trấn. Gió luồn qua những khe tường, làm rung lên những tấm biển hiệu han gỉ. Từng chữ cái rơi rụng, chỉ còn lại những mảnh vụn của ký ức. Một biển quảng cáo bằng thép cũ kỹ vẫn còn vẽ hình tàu, nhưng màu sơn đã bong tróc, chỉ còn lờ mờ dòng chữ cổ động: ``Thế giới vẫn đang chuyển động.''

Họ bắt đầu đi, từng bước rời khỏi nhà ga đã sụp đổ. Mỗi bước, họ đều cảm thấy mình đang mất phương hướng, không biết mình đang ở đâu, như thể đang lạc vào giữa mờ ảo và thực tại.

Con đường dẫn ra khỏi nhà ga như một dải lụa xám bạc, uốn quanh những căn nhà hai tầng cổ. Những cột đèn kiểu 1940 đứng lẫn giữa trụ đèn tín hiệu hiện đại năm 2018, một thứ chắp vá giữa hai thời kỳ. Dây điện chằng chịt vắt qua đầu, vài bóng đèn nhấp nháy yếu ớt dù không có nguồn điện rõ ràng. Trên mặt đường, vạch sơn phân làn vẫn còn nhưng bị nứt vụn, xen giữa là những tấm biển ``cấm vào'' nghiêng ngả, một vài biển vẫn mang biểu trưng của một công ty đã phá sản từ mấy chục năm trước.

\img{e3-1a}

Kuon đi đầu, ánh mắt lặng lẽ quan sát. Cậu chạm tay lên lan can sắt đã rỉ sét, từng lớp bụi dày phủ lên ngón tay. ``Cứ như\dots\ nơi này đã bị bỏ quên hàng chục năm vậy,'' cậu khẽ nói.

Kaigou đáp, giọng đều đều: ``\dots Em không nghĩ vậy. Có dấu của hai thời kỳ chồng lên nhau ở cả thời kỳ quá khứ và hiện đại này. Chắc chắn có một thế lực nào đó, hoặc là một cái gì đó\dots\ đã gộp chúng lại nơi này.''

Shino đi sau, bàn tay cô khẽ nắm vạt váy dài. Mái tóc nâu sẫm của cô khẽ lay trong gió, đôi mắt màu xanh biển long lanh như đang sáng lên trong khung cảnh âm u đáng sợ. Cô mím môi, dường như vẫn chưa chấp nhận hoàn toàn hình dạng của mình. ``Mình\dots\ vẫn không hiểu,'' cô nói khẽ, ``tại sao mình lại giống cô ấy đến thế\dots\ giống đến mức đáng sợ.''

Kuon quay lại nhìn cô, ánh mắt dịu đi. ``Có lẽ\dots\ bởi vì như Game-san nói đấy, cả hai là một phần của cùng một người,'' cậu đáp, ``chỉ là mỗi phần chọn cách đối mặt với sự việc khác nhau như thế nào.''

Nàng tóc nâu im lặng, không nói gì thêm. Cô nhớ lại lời giọng nói trầm mang sắc âm Pháp đó từng trấn an, rằng ``sự tồn tại của cô không phải sai lầm, mà là cần thiết để câu chuyện được trọn vẹn.'' Nghĩ tới đó, cô hít sâu, lấy lại bình tĩnh.

``\dots Ừm. Có lẽ là điều đó đấy.''

\ct

Con đường dẫn họ đến bờ sông Aogami, nơi nước lặng như gương. Bờ kè vẫn còn dấu vết của hàng lan can sơn đen, đã bong tróc và cong vênh.

\img{e3-1b}

Kaigou dừng lại, ánh mắt hơi xa xăm.
``Nơi này\dots'' cô khẽ nói, ``\dots là chỗ mà em - đúng hơn là Saikyou - từng đứng với Inari-san và Touka-san. Lúc đó trời cũng gần hoàng hôn thế này.''

Cô bất giác rùng mình, nhớ lại cái ánh mắt sắc lẹm của cô nàng tóc xanh ấy. Câu nói ``chỉ là cảm lạnh thôi'' đầy lạnh lùng và dứt khoát khiến cho cô vẫn không thể nào quên được, như một vòng lặp vọng lại từ ký ức xa xôi.

``\dots Mình cũng từng tới đây nữa,'' Shino bất chợt lên tiếng. Cả ba người ngoảnh mặt về phía cô, nghiêng đầu khó hiểu.

``\dots Mình cũng từng tới đây nữa,'' Shino bất chợt lên tiếng. Cả ba người ngoảnh mặt về phía cô, nghiêng đầu khó hiểu.

Cô bước lên bờ kè, ngón tay khẽ vuốt mảng hàng rào gỗ gãy gọn. ``Hôm đó trời cũng vàng úa như bây giờ\dots\ mình ngồi đây, cầm bông hoa cúc, nhìn bóng mình tan trong lòng sông.''

Shino khép mắt, giọng khẽ run. ``Mình vừa khóc vừa thì thầm: `Nếu có ai đó ở bên, mình sẽ không còn sợ nữa.' Rồi, ngày hôm sau mình lên ngôi đền, cầu mong Thần linh có thể giúp mình kết bạn\dots\ để không còn là cái bóng cô độc nữa.''

Cô mở mắt, nhìn phản ảnh mờ ảo trên mặt nước. ``Và có lẽ\dots\ từ khoảnh khắc mình gặp Reiko-san, mình đã hiểu rằng mình không còn là bản thân của trước kia nữa.''

Suzuki chỉ vào mình, tỏ vẻ hơi ngạc nhiên: ``M-Mình ư?''

``Là cậu đó, Suzuki-san. Chính xác hơn, là con người cũ của cậu,'' Shino gật đầu.

Kaigou khẽ đá chân lên viên gạch vỡ, tiếng vỡ lách tách. ``Nơi này như một bức tranh ghép: hai thời đại cố gắng hòa vào nhau nhưng chỉ tạo ra vết nứt.'' Cô nhìn Shino. ``Tôi thấy mình trong vết nứt đó.''

Kuon đưa tay gạt mái tóc bị gió thổi che mắt. ``Có vẻ như nơi này\dots\ cũng không phải là một nơi tầm thường như tôi tưởng.''

Suzuki ôm cánh tay mình, ánh mắt lướt qua những dây điện chằng chịt. ``Mình lại thấy\dots\ sự im lặng.'' Cô khẽ nở nụ cười gượng. ``Là thứ im lặng mà mình luôn tìm kiếm để không phải nghe tiếng tim mình đập. Nhưng giờ đây, tiếng đập ấy lại vang quá rõ\dots\ vì có ba người đứng cạnh.''

Cả bốn cùng nhìn xuống dòng sông phẳng lặng, bốn bóng hình in trên mặt nước, không còn là những cái bóng cô độc nữa, mà là một bức tranh ghép mới, chờ họ nối tiếp.

\ct

Họ tiếp tục lang thang trên thị trấn hoang tàn.

Khi qua khu phố truyền thống, sự pha trộn giữa cũ và mới trở nên rõ rệt hơn. Những ngôi nhà gỗ Nhật thời hậu chiến xen giữa tòa nhà bê tông kiểu hiện đại, biển hiệu vẽ tay nằm cạnh bảng đèn LED hỏng.

\img{e3-1c}

Suzuki chỉ tay về một ngã ba, chính xác là nơi mà cô cùng Kuon và Kaigou - đúng hơn, con người cũ của họ - tập trung: ``Đây là nơi chúng ta gặp Yuka-san, Akane-san, Miku-san, Saya-san và Kaoru-san\dots\ Lúc xuống tàu hôm đầu tiên.''

Cả bốn người dừng bước.

Con phố ấy giờ chỉ còn là tàn tích. Những tấm bạt quảng cáo rách nát bay phần phật, bên dưới là quán nước bỏ hoang, chiếc ghế gãy nằm chỏng chơ trước cửa. Tất cả gợi cảm giác như thị trấn bị giữ nguyên ở khoảnh khắc cuối cùng trước khi biến mất.

``Mình nhớ lúc đó mọi người còn rất vui,'' Suzuki khẽ nói, giọng hơi run. Cô ôm vai mình, mắt nhìn xuống vũng nước đọng trên nền đất. ``Giờ nhìn lại, cảm giác như đó chỉ là giấc mơ vậy.''

Cô chị bước lên, đôi mắt cô lướt qua từng mảnh vỡ. ``Không phải giấc mơ đâu.'' Cô cúi xuống nhặt một mảnh gỗ có dòng chữ ``nước đá chanh'' còn lem nhem. ``Là ký ức thật. Chỉ là\dots\ chúng bị bỏ lại.'' Cô siết chặt mảnh gỗ trong lòng bàn tay, như muốn giữ lại một phần của quá khứ. ``Chị từng nghĩ nếu quay lại đây, mọi thứ sẽ khác. Nhưng hóa ra, chỉ là bản thân chúng ta đã khác đi thôi.''

Cậu con trai đứng cạnh cửa hàng đóng cửa, nhìn vào tấm kính vỡ. Trong đó, phản chiếu bóng của bốn người họ - nhưng không phải bây giờ, mà là hình bóng của bản thân họ trước kia. ``Anh từng nghĩ nơi này là điểm bắt đầu cuộc sống mưới của chúng ta,'' cậu thì thầm. ``Nhưng giờ\dots\ sau những gì chúng ta đã trải qua, nó cũng có thể là điểm kết thúc.''

Cậu quay lại nhìn ba người kia, ánh mắt kiên định. ``Chúng ta không còn là những người đó nữa. Nhưng có lẽ\dots\ họ vẫn còn len lỏi đâu đó bên trong chúng ta.''

Shino lặng lẽ bước đến bên chiếc ghế gãy. Cô ngồi xuống, tay vuốt nhẹ lớp bụi. ``Mình chưa từng ở đây,'' cô nói, giọng nhỏ nhưng rõ ràng. ``Nhưng mình cảm được nỗi buồn của nơi này.'' Cô ngẩng lên, nhìn ba người kia. ``Các cậu đã từng có một nơi để quay về. Còn mình\dots\ mình chỉ có bóng tối.'' Cô mỉm cười, nhẹ nhưng không buồn. ``Nhưng bây giờ, mình có ba người đứng đây. Có lẽ, điều đó đủ để bắt đầu một nơi mới.''

Suzuki ngước nhìn Shino, mắt cô ánh lên niềm tin. ``Shino-san\dots\ mình không biết nói sao. Nhưng mình nghĩ, dù quá khứ có tan vỡ đi nữa, thì hiện tại vẫn còn đây.'' Cô bước lại, đứng bên cạnh cô bạn tóc nâu. ``Chúng ta có thể không còn là người cũ, nhưng chúng ta có thể là người mới - cùng nhau.''

Kaigou nhìn hai người, rồi quay sang Kuon. ``Anh nghĩ sao?''

Cậu không trả lời ngay. Cậu nhìn lên bầu trời xám xịt, rồi lại nhìn xuống con phố hoang. ``Anh nghĩ\dots'' cậu thở dài, nhưng không phải vì mệt mỏi. ``Anh nghĩ chúng ta đang đứng ở ngã ba. Không phải chỗ này,'' cậu chỉ xuống đất, ``mà là trong lòng.'' Cậy nhìn từng người. ``Chúng ta có thể chọn tiếc nuối, hoặc chọn bước tiếp. Chúng ta đã chọn bước tiếp - dù không biết đi đâu.''

Cả bốn im lặng. Gió thổi qua, làm tấm bạt quảng cáo rung lên như tiếng vỗ tay yếu ớt. Không ai nói thêm, nhưng cũng không ai rời đi. Họ chỉ đứng đó, giữa nơi từng là đầu và giờ là đống đổ nát - và trong lòng mỗi người, một hạt giống mới bắt đầu nảy mầm.

\ct

Họ rẽ qua vài con đường nhỏ hơn, nơi mặt đường nhựa đen đã lồi lõm thành những vảy rồng rắn, chỗ lộ ra lớp đá ong xám của thập niên bốn mươi. Mỗi bước đi là một tiếng vọng của quá khứ: bảng tên lớp học gỗ thông phủ sơn xanh đã bạc màu, chữ ``Trường Tiểu học Aogami - 1946'' còn đọng lại trong khe nứt; tờ rơi quảng cáo trường học đã phai chữ, chỉ còn dòng ``Khai giảng 2018'' loang lổ như giọt nước mắt in trên giấy; mảnh băng cassette cũ nằm lăn lóc trên vỉa hè, dải từ từ tuôn ra, phản chiếu ánh hoàng hôn như dải băng thời gian đứt đoạn.

Kuon cúi nhặt mảnh băng, ngón tay giật nhẹ khi lớp sắt gỉ đâm vào thịt. ``Nghe được gì không?'' cậu khẽ hỏi. Cả bốn im lặng, nhưng trong đầu họ dường như vang lên tiếng nhạc jazz cũ, tiếng cười học trò, tiếng loa phát thanh thông báo chiến tranh đã qua\dots\ tất cả hòa vào nhau thành một bản giao hưởng nghẹn ngào nhưng u ám tới mức đáng sơ.

Kaigou đưa tay vuốt vách tường đá ong, nơi lớp vữa đã rơi từng mảng, lộ ra lớp gạch ta đỏ au của những năm tháng đói kém. ``Như thể ai đó cầm cưa cắt ngang dòng thời gian rồi dán hai đầu lại bằng băng keo,'' cô thì thầm. ``Chúng ta đang bước trên chỗ keo đó, chỉ cần giẫm mạnh một chút, mọi thứ sẽ rách toạc.''

Shino khựng lại trước một cửa hiệu may cũ. Trên tấm kính vỡ, ma-nơ-canh thời chiến mặc áo dài tay, bên cạnh là chiếc váy ngắn sequin của năm 2017, cả hai cùng phủ bụi như những xác ướp bị quên lãng. Cô đưa ngón tay viết một dấu hỏi trên lớp bụi, rồi lặng lẽ xóa đi, như sợ câu hỏi về năm tháng sẽ vang vọng mãi không dứt.

Suzuki cúi xuống nhặt chiếc huy hiệu - dường như là lễ hội văn hóa nào đó, mặt sau vẫn còn khắc tên cô một ai đó đã mờ đến mức chỉ còn là những nét gợn. Cô siết chặt trong lòng bàn tay, móng tay đâm vào thịt. ``Mình từng nghĩ 2018 là năm của riêng chúng ta,'' cô nói, giọng khản đặc. ``Nhưng hóa ra nó chỉ là một lớp sơn phủ lên 1946 - và cả hai đều đang bong tróc.''

Thời gian như tan chảy, không còn biết đâu là 1946, đâu là 2018. Một cơn gió thổi qua, cuốn theo lá khô của hai mùa thu cách nhau bảy mươi hai năm, xoáy thành một cơn lốc nhỏ trước mặt họ. Trong mắt mỗi người, những vết nứt của thời gian hiện lên như những vết rạn trên gương - không đau, nhưng phản chiếu muôn ngàn mảnh hình họ từng là, đang là, và có thể sẽ không bao giờ thành.

\ct

Họ rẽ khỏi con phố chính, băng qua một con dốc đất đỏ trơn trượt, hai hàng cây bạch đàn gai góc vươn cành khô ra như những bàn tay xin nước. Gió từ rừng thông thổi xuống, mang theo mùi nhựa và lá mục, xoáy vào những vết thương hở của thị trấn như muốn rút hết hơi ấm còn sót lại. Mỗi bước chân nện lên mặt đường đều kéo theo tiếng sỏi lăn lách cách, vang lên giữa khoảng không tĩnh lặng đến nghẹt thở.

Kuon dừng lại trước ngã ba có tấm biển gỗ mục nát chỉ còn dòng chữ ``Trường Trung học Aogami - 1,2 km'' in lờ mờ bằng sơn trắng. Cậu đưa tay gạt lớp rêu bám trên mũi tên, lớp sơn vụn ra như vôi. ``Tôi vẫn nhớ lần đầu tiên khi cùng Kaigou-chan đến trường cùng với Suzuki.''

Cô em út rụt rè, gật đầu: ``Em cũng có nhớ\dots\ Lần đầu tiên em đến đó để nhận bộ dồng phục trường\dots\ và cũng là lần đầu em gặp mặt bọn họ. Nanase-san, Yae-san, Sakura-san, Shiina-san và Honoka-san ở đó.''

Con đường mòn nhỏ hẹp dần, lát đá ong lộ ra từng khối vuông vỡ góc, xen lẫn rễ cây nổi lên như những con rắn gỗ đang căng cứng. Hai bên, bụi tre già lay động ken két, để lộ những mảnh băng cassette vương vãi, dường như ai đó đã từng ngồi đây, thu âm lại tiếng gió rừng làm nhạc nền cho một bài hát không bao giờ kịp phát hành.

Kaigou cúi nhặt một cuộn băng, dải từ đã rối thành mớ bòng bong đen xì. Cô giật nhẹ, tiếng từ xé gió nghe như tiếng khóc thương. ``Nghe giống như thứ âm thanh mà Saikyou từng ghét,'' cô thì thầm, ``lúc còn cố gắng làm một học sinh bình thường.''

Suzuki ôm sát chiếc áo khoác, mắt lia qua những tàn lều lá dựng bằng cành cây và ni lông, có lẽ là nơi trú ẩn tạm bợ của nhóm học sinh nào đó trong những ngày tháo chạy khỏi vòng lặp. Cô khựng lại khi thấy một đôi giày thể thao cũ nằm lẫn trong đám lá khô, dây giày thắt chặt như thể chủ nhân vẫn còn đó, chỉ vừa biến mất trong chớp mắt. ``Em\dots\ em không dám chắc mình muốn nhìn thấy những gì ở trường đâu,'' cô nói, giọng run nhẹ.

Shino bước tới bên cô, đặt tay lên vai Suzuki. ``Chúng ta đã bước qua chỗ sợ hãi rồi,'' cô khẽ đáp, ``dù kết quả có là gì, ít ra chúng ta cũng không còn phải đối mặt một mình nữa.'' Ánh mắt cô hướng về phía trước, nơi tán thông già chắn ngang lối đi, tạo thành một vòm tối như miệng hang. Bên kia vòm cây, ánh hoàng hôn đổ xuống thành những vệt vàng cam loang lổ trên mái tôn dập sóng của một tòa nhà cao tầng.

\ct

Họ bước qua vòm thông. Gió đột ngột ngừng thổi, như thể khu rừng cũng nín thờ chờ đợi. Trước mặt là Trường Trung học Aogami.

Suzuki khẽ nuốt nước bọt. ``\dots Chúng ta thật sự quay lại rồi.''

Kaigou nheo mắt nhìn. ``Hoặc là, chúng ta chưa bao giờ rời khỏi đây.''

\img{e3-1d}

Cánh cổng sắt bị móp méo, lớp sơn gỗ bị bong tróc từng mảng lớn, để lộ những đường màu nâu sần sùi như mạch máu khô. Bảng tên trường nghiêng lệch, chỉ còn ba chữ ``Trung học Aogami'' lem nhem. Những ô cửa kính vỡ toang như mắt người khuyết, rèm rách tả tơi bay phần phật trong cơn gió thoáng qua, phía sau là bóng tối nuốt chửng hành lang dài hun hút. Trên tường, rêu phong bò lên gần hết tầng hai, xen lẫn những dòng graffiti cũ kỹ: ``Chúng tôi vẫn ở đây'' chữ ``ở'' bị cào xước nguệch ngoạc, như thể ai đó đã định xóa đi rồi lại thôi.

Kuon đẩy nhẹ cánh cổng. Tiếng kim loại ken két vang lên, gợn lên những đợt âm thanh rít xé không khí. Một con quạ đen đột ngột bay vút ra từ mái hiên, vỗ cánh lạch xạch như tiếng vỗ tay chế giễu. Bốn người đứng im, nhìn theo bóng chim khuất sau rặng thông, rồi cùng quay lại nhìn ngôi trường, nơi từng là điểm tựa của tuổi trưởng thành, giờ chỉ còn là xác xương bê tông giữa bầu trời xám xịt.

Shino đứng sau, nét mặt hiện vẻ do dự nhưng cũng không giấu được sự tò mò: ``Bên trong\dots\ có gì chờ đợi chúng ta sao?''

Cậu con trai nhìn cô, rồi nhìn hai người còn lại.
Ánh mắt anh ánh lên sự kiên định quen thuộc. ``Chỉ có một cách để biết thôi.''

Cả bốn người đứng trước cánh cổng gỗ mở hé, ánh sáng mờ mịt từ đám mây phủ lên ngôi trường hoang tàn một lớp sắc xám xịt lạnh lẽo.

\ct

Bên trong là khu để giày cũ kỹ, nơi từng chào đón hàng trăm học sinh mỗi sáng. Giờ đây, những chiếc tủ gỗ đã mục nát, sơn bong tróc, khóa sắt hoen gỉ. Hầu hết nhãn tên trên ô tủ đã bị xé hoặc mờ, chỉ còn vài dòng chữ đứt quãng: nét bút mờ nhòe, vài cái bị rách, vài cái bị dán đè lên bởi giấy cũ. Một cơn gió nhẹ lùa qua, làm vài tờ nhãn rời bay lả tả xuống nền gỗ đã loang lổ.

\img{e3-1e}

Suzuki cúi xuống, nhặt một mảnh giấy còn đọc được vài nét mực: ``Shi\dots\ Yae.'' Cô khẽ nhíu mày.

``Sao vậy?'' chị lớn của cô hỏi.

``Đó chẳng phải\dots\ là một trong những học sinh cũ của nơi này sao?'' cô ấy lẩm bẩm, gần như thì thầm. ``Họ đều là những người\dots\ đã chết trong vụ sạt lở đất năm đó.''

Shino đứng bên, mặt thoáng biến sắc.

Khoảnh khắc mà cô bé mắt hai màu nhắc tới họ khiến cô nhớ đến những tiếng cười, những lời chọc ghẹo trong lớp, và cả buổi tối sau trận mưa hôm ấy, khi họ lẻn vào trong trường để khám phá.

``Là lỗi của mình\dots'' cô khẽ nói, giọng nghẹn lại. ``Nếu hôm đó mình can đảm hơn, có lẽ họ đã không--''

Kuon vội cắt ngang, đặt tay lên vai cô: ``Không ai có thể biết trước chuyện gì sẽ xảy ra, Shino-san. Chuyện đã qua rồi thì cứ để qua thôi. Giờ ta nên hiểu rõ vì sao mọi thứ lại trở thành như thế.''

Kaigou bước tới, ngồi xuống bên cạnh một tủ giày đã sụp nửa phần. Trên cánh cửa còn mảnh giấy bị xé dở, chỉ còn lại một chữ: ``Rei--''.

Cô thoáng im lặng. ``\dots Reiko.''

Suzuki giật mình, mắt mở to. ``Reiko\dots\ đó là em, đúng không, Onee-sama?''

Cô chị không đáp ngay, chỉ nhìn cô một lúc lâu rồi gật nhẹ. ``\dots Chính xác hơn, là con người cũ của em đấy.''

``Đ-Đúng là như vậy\dots\ Kỳ lạ thật đấy. Tại sao em vẫn còn nhớ những ký ức từ con người cũ của mình nhỉ\dots''

Câu nói đó vang lên trong không gian trống, nặng nề đến nghẹt thở.

Kuon nhìn chiếc đồng hồ đeo tay, màn hình trên mặt đồng hồ chỉ còn một vạch pin yếu ớt. Cậu khẽ gõ nhẹ lên mặt đồng hồ, như thể đang gọi ai đó. ``Game-san, ông còn đấy chứ?''

Một tiếng nói nhỏ, trầm, pha chút âm hưởng Pháp vang lên từ chiếc đồng hồ: ``Tôi vẫn ở đây\dots\ nhưng pin thì không còn nhiều đâu. Có gì thì nói nhanh đi.''

``Ông có thể giải thích cho chúng tôi không? Tại sao chúng tôi vẫn còn giữ ký ức từ những con người cũ của mình?''

``Đó là vì,'' giọng nói đáp lại, có chút rệu rã nhưng vẫn có gì đó rõ ràng, ``những ký ức ấy không thuộc về 'ai đó' khác. Chúng chỉ là phiên bản khác thôi. Khi vòng lặp vỡ, các cậu không `mất đi' một phần, mà chúng \emph{hợp nhất} lại với nhau. Những gì các cô cậu nhớ là những gì mà con người trước đó của các cậu đã từng trải qua.''

``Ra thế,'' Kaigou đáp. ``Bảo sao mà chúng tôi vẫn còn nhớ được những sự kiện ở đây, mặc dù chúng tôi hồi đó đã ở trong một con người khác.''

``Chúng có thể hữu ích cho việc phá giải những bí ẩn ở nơi này đấy,'' ông ta tiếp tục. ``Hãy tận dụng thật tốt.''

``Vậy chúng cháu nên bắt đầu từ đâu ạ?'' Suzuki băn khoăn.

``\dots Sao không bắt đầu từ nơi mà căn nguyên mọi chuyện xảy ra nhỉ? Đó là bước đầu để tìm ra lời giải đấy.''

Cậu con trai chau mày, ngẫm nghĩ điều mà trợ lý ảo của cậu vừa nêu.  ``Ông nói cũng có lý.''

Shino đứng im, lặng lẽ nhìn ba người kia. Gió thổi qua hành lang, làm tung bay mái tóc nâu sẫm của cô. ``Dù rằng ông ấy hơi kỳ quặc,'' cô khẽ nói, ``nhưng ông ấy cũng để lại một lời giải thích để chúng ta tự tìm câu trả lời cho riêng mình.''

Cô nhìn cả ba, chậm rãi hít một hơi.  ``Có lẽ chúng ta nên\dots\ đến phòng Mỹ thuật. Mình nhớ\dots\ nơi đó có thứ mà Honoka-san từng chia sẻ. Cũng là nơi bắt đầu của lời đồn ngày ấy.''

Kuon nghiêng đầu. ``Lời đồn?''

Cô gật đầu, ánh mắt hơi xa xăm: ``Hồi đó, cô ấy bày ra một trò thử thách lòng can đảm. Cô ấy khăng khăng muốn biết bên trong phòng Mỹ thuật đó có cái gì. Nếu mình không nhớ nhầm, thì đó là\dots\ hồn ma của cô gái đã dùng máu của mình để vẽ bức tranh thì phải.''

Cả nhóm im lặng. Tiếng gió từ hành lang sau lưng rít qua những khe cửa, như vọng lại tiếng thì thầm cũ kỹ của những học sinh đã mất.

Họ bắt đầu đi dọc theo hành lang chính.

\img{e3-1f}

Trần nhà nứt nẻ, từng mảng bê tông tróc ra, để lộ những cục đất và đá vụn còn sót lại từ trận sạt lở định mệnh. Mỗi bước đi phát ra tiếng vang khô khốc, vang vọng trong không gian hun hút như tiếng bước chân của quá khứ đang bám theo họ. Ánh sáng yếu ớt lọt qua những khe cửa sổ vỡ, làm nổi bật những dòng chữ cũ kỹ đầy sắc đỏ, như thể viết bằng máu: ``Chúng tôi vẫn ở đây'', ``Lớp 1-C không bao giờ tan'', ``Ai đó hãy nhớ tên tôi''\dots những dòng chữ nguệch ngoạc như tiếng kêu cứu bị thời gian bỏ quên. Trong con mắt của bọn họ, nơi này thậm chí còn xập xệ và hỗn loạn hơn cả những gì mà họ còn nhớ.

Dưới chân, những mảnh kính vỡ phản chiếu ánh hoàng hôn nhợt nhạt, xen lẫn những mẩu báo úa vàng rơi rải rác, vài tờ còn đọc được dòng tít mờ mờ: ``Thảm kịch tại Aogami: năm học không bao giờ kết thúc.'' Một chiếc giày học sinh cũ nằm lẻ loi dưới gầm tủ, dây buộc vẫn còn thắt chặt, như thể chủ nhân chỉ vừa biến mất trong chớp mắt. Trên tường, những bức ảnh lớp học bị rách một nửa, khuôn mặt học sinh bị cào xước đến mức không còn nhận ra ai là ai, chỉ còn lại những nụ cười gượng gạo, như đang cố níu giữ một ký ức không còn nguyên vẹn.

Suzuki bước khẽ, tay nắm chặt vạt váy, mắt không rời những dấu vết trước mặt. Cô cố gắng không nhìn vào chiếc áo khoác đồng phục còn vương vãi trên sàn - có lẽ là của một học sinh năm ba - với dòng chữ thêu tên của một người đã bị máu nhuốm đen, cứng đặc. Mỗi bước chân của cô như đang dẫm lên một phần ký ức mà cô từng cố gắng chôn vùi.

Kaigou đi sau, lặng lẽ quan sát từng chi tiết, đôi mắt cô ánh lên một nỗi buồn khó gọi tên. Cô dừng lại trước một bảng tin cũ, nơi những tờ thông báo về lễ hội văn hóa, buổi họp phụ huynh, và cuộc thi hát vẫn còn dán lệch lạc, như thể thời gian đã ngừng trôi từ ngày đó. Một tờ giấy nhỏ còn sót lại ghi: ``Em sẽ đợi anh ở cầu thang sau giờ học - XX/XX/2018.'' Cô khẽ chạm vào dòng chữ, như thể muốn cảm nhận hơi ấm của người đã viết ra nó, nhưng chỉ còn lại một sự lạnh lẽo.

Còn Kuon ở đằng trước, bước đều và chắc, nhưng bàn tay cậu khẽ co lại, như thể cũng đang kìm nén điều gì đó. Cậu không dám nhìn vào lớp học cuối hành lang, nơi cửa vẫn hé mở, bên trong bàn ghế ngổn ngang, và trên bảng vẫn còn dòng chữ phấn: ``Bài kiểm tra ngày mai: Chương 3, 4, 5.'' Một chiếc cặp học sinh vẫn còn để trên bàn thứ hai từ dưới lên, như thể chủ nhân chỉ vừa ra ngoài một lát, nhưng bụi đã phủ dày, và cánh cửa sổ đập vào khung kêu ken két theo nhịp gió.

Cậu bất giác dừng bước, ánh mắt cậu mắc kẹt trên vết nứt hình chân chim chạy dọc tường lớp cuối hành lang. Trong khoảnh khắc, tiếng gió thay bằng tiếng rú của bê tông đang gãy, tiếng kính vỡ như mưa. Cậu thấy mình không còn đứng ở ngôi trường lúc này nữa nữa; mà là Kyuen trong đống đổ nát của chính nơi này thời hiện đại, nơi dầm sắt xiên qua bục giảng, nơi mái ngói sập xuống bẹp dí hàng ghế màu xanh lá. Mùi bụi xi-măng, mùi máu tanh, mùi khói báo động hòa vào nhau. Cậu nhớ cô em gái của cậu đã khản cổ gọi những người bạn hồi đó, nhưng chỉ nhận lại tiếng vang rỗng tuếch từ khối bê tông đang chúi dần xuống.

Cậu nắm chặt tay, móng đâm vào lòng bàn tay, đau đến mức buộc mình trở về hiện tại. Cậu thì thầm, giọng khàn đặc: ``\dots Lại là ký ức đó à.''

Kaigou đứng ngay sau lưng cậu, cũng bị bất động. Đôi mắt cô không còn nhìn vào bảng tin nữa; mà nhìn xuyên qua nó, thấy Saikyou đang bới tung đống gạch vỡ tìm bất kỳ dấu hiệu sống sót nào. Cô nhớ cái cách mà người anh trai song sinh của cô đã kéo cô ra khỏi đống đổ nát, nhớ tiếng khóc nghẹn khi cô nhận ra chỉ còn hai anh em đứng thở. Cô đưa tay lên vuốt nhẹ vết sẹo cũ trên cánh tay, vết sẹo mà ở dòng thời gian này chưa từng tồn tại, như thể có thể gợi lại cảm giác đau đớn để chứng minh rằng những người bạn ấy đã từng sống.

Một giọt nước mắt rơi xuống nền gạch nứt, tan biến vào lớp bụi phủ. Cô không khóc thành tiếng; chỉ thì thầm: ``Khoảnh khắc đó\dots\ đúng là nghiệt ngã thật.''

Cả hai đứng song song, cùng nhìn vào căn phòng học bỏ hoang, nhưng thực ra đang nhìn vào một ngôi mộ chung không bia, không tên, không tuổi. Gió thổi qua hành lang, cuốn theo tro bụi của hai thời đại, thì thầm như an ủi: ít ra, lần này họ không còn đơn độc khi nhớ về họ.

Shino đi cuối cùng, mắt không ngừng đảo quanh, mọi ngóc ngách đều khiến cô thấy quen thuộc đến lạnh người. Cô dừng lại trước một chiếc gương vỡ nứt dọc hành lang, bề mặt phản chiếu khuôn mặt cô méo mó, chia cắt thành những mảnh vụn. Trong khoảnh khắc, cô tưởng như nhìn thấy một cô gái khác, cũng tóc nâu, nhưng có đôi mắt nâu đang dạo bước cùng với cô, cảm tưởng như bản thân vẫn đang ở trong một cơ thể nào đó khác. Cô thụt lùi một bước, tim đập thình thịch, nhưng sau lưng chỉ còn bóng tối và tiếng gió rít qua khe cửa.

Một phần ký ức, một phần ám ảnh.

Thỉnh thoảng, cô còn tưởng như nghe thấy tiếng cười của ai đó vọng ra từ tầng hai, rồi tan biến vào hư vô, như tiếng chuông báo vào lớp vang lên một lần cuối, rồi im bặt mãi mãi.

\ct

Cuối cùng, họ dừng lại trước phòng Mỹ thuật.

Cánh cửa gỗ sậm màu đã ngả sang xám tro, bảng tên trên cao thì rách rưới, nhưng vẫn đủ để mọi người có thể đọc được tên của căn phòng. Trên bề mặt gỗ, những vết xước hình như là dấu móng tay, hay có thể là lưỡi dao, tạo thành những đường xiêu vẹo nhỏ như mạng nhện. Một mảnh băng dính cũ còn dính lại ở góc cửa, phai màu thành vàng úa, in hình một con bướm đã bị rách đôi.

Cánh cửa mở hé, rèm bên trong khẽ lay, mùi bụi và sơn khô thoang thoảng. Ánh sáng cuối ngày lọt qua khe cửa, chiếu lên một vệt bụi bặm trên nền gạch, nơi dường như từng có ai đó đứng lâu đến mức in hình bàn chân. Một chiếc ghế nhỏ gỗ lim bị kéo lệch ra khỏi tường, trên mặt ghế còn vết sơn dầu loang lổ, có lẽ là dấu vết của một bức tranh dở dang.

``Chúng ta quay lại nơi này rồi,'' Kuon khẽ cất tiếng, đôi mắt dán vào miếng dán phong ấn đã bong tróc.

Suzuki dừng bước, liếc nhìn hai người đi sau rồi quay sang Shino. ``Đây là nơi cậu nói đến sao, Shino-san?''

Shino khẽ gật đầu. ``Phải\dots\ đó chính là nơi mà bọn mình đã tham gia cuộc thử thách đó. Mình với Kyuen, Saikyou và các bạn học khác đã lục tung cả căn phòng đó để tìm vết tích về lời đồn đó, nhưng\dots''

``\dots không tìm thấy sao?'' Kaigou tiếp lời. ``Bản thân tôi cũng nhớ rõ mà.''

Suzuki ôm cánh tay mình, mắt nhìn vào bóng tối bên trong. ``Nhưng lần này\dots\ có lẽ chúng ta khác rồi. Có thể\dots\ có thể chúng ta sẽ thấy thứ mà trước đây mình không thể thấy.''

Cả bốn người đứng lặng, mắt không rời khe cửa hé mở như thể đang chờ một lời mời từ quá khứ.

``\dots Chỉ còn một cách để biết thôi,'' cô nàng tóc nâu đáp.

Cậu con trai siết chặt tay, đốt ngón tay lạnh ngắt. ``Nếu có gì đó còn sót lại, nó sẽ phải đối mặt với chúng ta.''

Cô nàng tóc đuôi ngựa gạt mái tóc bị gió thổi che mắt, ánh nhìn sắc như dao cùn. ``Dù là bóng ma hay ký ức, tôi sẽ kéo nó ra ánh sáng.''

Suzuki hít một hơi thật sâu, môi mím lại thành một đường thẳng. ``Chúng ta không còn là những đứa trẻ chạy trốn nữa.''

Bốn bóng hình đứng thành một hàng, bước vào phòng Mỹ thuật như bốn mũi tên cùng được nockẻn trên một dây.

Tiếng gió rít qua khe cửa, mang theo mùi mốc ẩm và nỗi buồn của quá khứ, như thể ngôi trường vẫn còn nhớ những linh hồn đã rời bỏ nó.

Kuon nhẹ đẩy cánh cửa phòng Mỹ thuật, tiếng bản lề gỉ sét rít lên một cách khó chịu, như thể cánh cửa phản đối việc bị mở ra sau khoảng thời gian dài đằng đẵng kia.

*KẸẸẸẸẸẸẸẸẸẸT*

Ngay khi bước vào, một cơn đau đầu nhói buốt đồng loạt ập đến.

Kaigou khụy nhẹ xuống, một tay chống vào tường, mắt nhắm nghiền lại. ``G-Gì thế này\dots?'' cô rên khẽ, trán đổ mồ hôi lạnh.

Shino thì tệ hơn. Cô loạng choạng vài bước, tay ôm đầu, hơi thở đứt quãng. ``Đau\dots\ như ai đó đang\dots\ đâm vào bên trong trí nhớ của mình vậy\dots''

Suzuki vội đỡ lấy cô, giọng lo lắng: ``Shino, cậu ổn chứ?!''

Kuon liếc quanh, nheo mắt chịu đựng cơn nhói trong thái dương.
``Tại sao\dots\ bên trong lại bừa bộn tới vậy\dots?'' giọng cậu lạc hẳn đi.

Căn phòng trước mắt họ thậm chí còn tồi tệ hơn những gì họ thấy trong giấc mơ, hay cả trong những vòng lặp trước đó, như thể đây mới là ``phiên bản nguyên thủy'' mà các ký ức đã cố tình che giấu.

Những giá vẽ ngổn ngang dựng nghiêng, vài chiếc đổ sập xuống sàn, khung gỗ gãy đôi. Ghế đẩu vương vãi khắp nơi, nhiều chiếc gãy chân, số khác phủ kín bụi.

Trên tường, các bức tranh treo san sát nhau -  phong cảnh, chân dung, tĩnh vật - giờ đã bị xé, nhòe, hoặc rách vụn, sơn khô bong thành từng mảng. Nhiều tấm tranh rủ xuống, để lộ lớp tường phía sau nứt nẻ và đen ám bụi.

Góc phòng, tủ đựng dụng cụ lật nghiêng, hộp phấn vỡ vụn, lọ sơn chảy khô thành từng vệt sẫm như máu.
Mùi sơn dầu, mốc và gỗ mục hòa vào nhau tạo thành một thứ hương khó chịu, nặng nề, như đang quẩn quanh trong không khí không lối thoát.

\img{e3-1g}

Suzuki tiến lên trước vài bước, đôi giày chạm vào mảnh kính vỡ phát ra tiếng ``rắc'' nhỏ. Cô ngẩng đầu lên, rồi khựng lại.

``K-Kia chẳng phải là\dots''

Ở chính giữa căn phòng, chỉ còn lại một bức tranh nguyên vẹn. Nó được treo trên giá vẽ duy nhất còn đứng thẳng, xung quanh là vô số mảnh vụn tranh cũ.

Cả bốn người đồng loạt bước tới, và cùng lúc, hơi thở của họ nghẹn lại.

Trong tranh là hình ảnh một cô gái trẻ, đôi mắt mở trừng trừng nhìn thẳng ra phía họ. Những nét chì đen đan xen với những vệt màu đỏ tươi như máu rỏ giọt.  Trên môi cô là biểu cảm méo mó, nửa như một nụ cười, nửa như một tiếng thét bị đóng băng.

\img{e3-1h}

\dots

\dots

``\textbf{Gyaah!}''

Shino chợt giật mình thét khẽ, ngã ngửa ra sau, đập vào giá vẽ khác khiến chúng đổ ầm xuống.

``\dots Cái quái?'' Game bất chợt thốt lên, giọng hiện rõ sự khó tin.

Kaigou và Suzuki theo phản xạ đưa tay che mắt, còn Kuon thì sững người, trán nhăn lại, ánh mắt đờ ra trong vài giây, trông như cậu vừa nhìn thấy thứ gì vượt quá ngưỡng hiểu biết của con người.

``B-Bức tranh quái quỷ gì đây\dots?'' cô nàng tóc nâu lắp bắp, bàn tay run rẩy chỉ vào bức họa. ``Tại sao trông nó lại\dots\ như đang nhìn thẳng vào chúng ta\dots?''

Không một ai đáp. Không đơn thuần chỉ là vì họ không có lời giải, mà họ cũng bị choáng với bức họa đó.

Chỉ có tiếng gió luồn qua khe cửa, hòa cùng tiếng nhỏ giọt lách tách từ trần nhà, nghe như thể căn phòng đang thở.

``\dots Hình như tôi vừa nhớ ra cái gì đó rồi. Chính là nó.''

Giọng của cô chị gái cất lên, khiến ba người còn lại ngoái nhìn cô. Cô hít sâu, khẽ nói, giọng chậm rãi: ``Bức tranh đó\dots\ là lời nguyền' trong phòng Mỹ thuật của trường Aogami. Cô gái đã dùng chính máu của mình để vẽ nên một kiệt tác, cho đến khi cô ta chết vì mất máu. Và kể từ đó, bất kỳ ai nhìn vào bức tranh đều\dots''

``\dots bị nguyền rủa tới chết,'' Suzuki nói tiếp, giọng nhỏ dần.

Shino lùi lại một bước, ánh mắt vẫn không dứt khỏi bức tranh. ``Không\dots\ không thể nào\dots\ nó chỉ là một lời đồn học sinh thôi mà.''

Kuon khẽ lắc đầu. ``Không đâu. Đây chính là thứ mà chúng ta từng thấy trong vòng lặp đầu tiên. Nhưng\dots\ lúc đó, bức tranh này không hề thật đến mức này.''

Một khoảng lặng kéo dài.

Không ai trong bốn người biết nên làm gì tiếp theo. Cảm giác như căn phòng đang nhìn họ, như chính bức tranh đang thở, mời gọi họ tiến gần thêm.

``Có lẽ\dots\ bức tranh này đang có gì đó mời gọi chúng ta,'' Game đưa ra gợi ý. ``Nó dường như muốn nói với ta một điều gì đó.''

``Ông\dots\ ông chắc chứ?'' cô gái tóc nâu lắp bắp hỏi, mắt đưa qua đưa lại thể hiện chần chừ.

``Khả năng cao. Dường như nó cũng không có ý gì muốn làm hại chúng ta, ngoài việc trông nó kinh khủng và tởm lợm ra cả.''

Game nói xong, không khí như nặng thêm một chút. Bốn người đứng đó, lặng im, mắt không rời bức tranh - thứ duy nhất còn sống sót giữa đống đổ nát. Có điều gì đó trong ánh nhìn của họ, không phải sợ hãi, mà là một thứ gọi mời kỳ lạ, như thể bức tranh đang thì thầm vào tai họ: ``Nếu các người dám, thì hãy đến đây.''

Không ai bảo ai, Kuon, Kaigou, Suzuki và Shino cùng tiến lại, giơ tay chạm vào khung tranh.

Ngay khoảnh khắc đầu ngón tay chạm vào mặt sơn khô lạnh lẽo, cả căn phòng rung lên.

Tiếng gió bị nuốt mất, ánh sáng ngoài cửa sổ vụt tắt, thay vào đó là một luồng sáng đỏ thẫm như máu tràn ngập không gian.

Shino kêu khẽ, cảm giác như rơi vào một khoảng không vô tận.

Kaigou ôm lấy đầu, còn Suzuki nắm chặt tay Kuon, nhưng cả hai cũng nhanh chóng mất thăng bằng.

Tiếng tim đập vang lên dồn dập, rồi im bặt.

Trong khoảnh khắc cuối cùng trước khi mọi thứ biến mất, bức tranh khẽ đổi hình, nụ cười méo mó kia như cong rộng thêm, thành một biểu cảm vừa đau đớn vừa mãn nguyện.

Âm thanh như pha lê vỡ vang lên. Không gian xung quanh họ nứt ra, những đường rạn trắng xóa lan tỏa khắp bức tường, khắp mặt đất, khắp không khí. Ánh sáng trắng đặc quánh tràn qua những khe nứt, nuốt lấy tất cả trong một chớp mắt.

Shino đưa tay lên che mắt. Trong giây lát, mọi thứ xung quanh họ tan chảy như sáp, rồi hòa vào nhau, thành một cơn lốc ánh sáng đang xoáy tròn.

Cảm giác như thể chính cơ thể họ bị bóp méo, kéo căng ra, rồi cuốn sâu vào trong một cơn xoáy vô hình.

``Chuyện\dots\ chuyện gì đang xảy ra vậy!?'' cô kêu lên, giọng cô vang vọng trong không trung méo mó.

Kaigou nghiến răng, cố giữ thăng bằng: ``Không phải dịch chuyển bình thường\dots\ giống như bị hút vào chính thứ gì đó.''

Rồi, tất cả tan biến, chỉ còn lại một màu trắng xóa, trắng đến mức không còn chút bóng đổ\dots

\il

Nhiều năm về trước đó.

Ánh hoàng hôn len qua khung cửa sổ, loang loáng hắt lên những mảng tường loang lổ, chiếu lên cô gái trẻ đang ngồi trước giá vẽ. Mái tóc đỏ của cô hơi rối, đôi mắt hai màu chăm chú dõi theo từng nét bút run run trên tấm toan trắng.

\img{e3-1i}

``Hừm\dots'' cô lẩm bẩm, giọng khẽ trôi giữa không gian tĩnh lặng. ``Dù có vẽ thế nào\dots\ nó vẫn không ra được như mình tưởng tượng.''

\chrint

\begin{wrapfigure}[14]{l}{6.5cm}
  \includegraphics[height=8cm]{../../../../Image/Book?/mari(n).png}
\end{wrapfigure}

Đó là Akagane Mari (\ruby{赤}{あか}\ruby{鐘}{がね}\ruby{瑪}{ま}\ruby{莉}{り}, \ruby{A-ca}{Hồng}-\ruby{ga-nê}{Chung}\ \ruby{Ma}{Mã}-\ruby{ri}{Lị}), cô gái có niềm đam mê mãnh liệt với mỹ thuật.

Dù dáng vẻ có phần trầm lắng, cô lại dễ gần và luôn biết cách khiến người đối diện cảm thấy ấm áp. Mỗi chiều, khi hành lang lớp học vắng người, Mari lại lặng lẽ mang bảng màu vào phòng mỹ thuật, miệt mài pha sắc, đặt xuống từng nét vẽ như thể đang gieo một phần tâm hồn mình lên đó.

Ước mơ duy nhất của cô là tạo ra một kiệt tác không ai có thể bắt chước, bức tranh khiến người xem phải dừng bước và không thể rời mắt.

Thế nhưng, đôi khi cô lại ôm đầu như thể một cơn đau đang búng ra từ sâu trong hộp sọ, và sau những lần ấy, ánh mắt cô như xa xăm hơn, như thể đã trả một phần của chính mình để đổi lấy nét vẽ hoàn hảo kia, cái giá mà chỉ riêng cô mới hiểu rằng mình đang chuẩn bị đón nhận.

\chrint

*XOẠCH*

Tiếng cửa trượt khẽ mở.

Một cô gái khác bước vào, đồng phục học sinh đã được xắn gọn, cổ tay đeo một chiếc vòng đỏ thẫm với ba chiếc chuông nhỏ leng keng khe khẽ theo từng bước chân.

``Cậu đây rồi, Mari-san,'' cô cất giọng, nhẹ nhưng vang lên rõ trong căn phòng yên tĩnh. ``Cậu làm gì ở đây vào giờ này thế? Hết giờ học rồi còn gì\dots''

Cô nhìn vào giá tranh mà Mari đang dựng. Nét lo lắng hiện rõ trên gương mặt của cô:``Vẫn còn chăm chú vào bức tranh đó sao?''

Mari không quay lại. Cô chỉ ``ừm'' một tiếng như để xác nhận, rồi đặt cọ xuống, ngước lên nhìn bức tranh dang dở trước mặt.

``Reiko-san, mình muốn vẽ một bức tranh mà không ai khác có thể vẽ được. Một thứ\dots\ thật sự đặc biệt.''

Reiko hơi nghiêng đầu. ``Bức tranh mà không ai có thể vẽ được?''

``Phải.'' Cô nàng tóc đỏ đáp, ánh mắt cô phản chiếu sắc đỏ của hoàng hôn như đang rực cháy. ``Một bức tranh khiến người ta phải dừng lại\dots\ không phải vì nó đẹp, mà vì họ không thể rời mắt đi được. Thứ gì đó\dots\ chạm tới tận nơi sâu nhất của trái tim.''

Reiko chớp mắt, thoáng bối rối: ``Nghe có vẻ hơi đáng sợ đấy, Mari-san.''

Mari bật cười khẽ, nhưng trong nụ cười ấy có chút gì không ổn. ``Có lẽ vậy. Nhưng nghệ thuật mà. Đâu phải lúc nào cũng dễ chịu. Đôi khi để tạo nên `kiệt tác', cậu phải trả giá một phần nào đó của chính mình.''

Cô nàng tóc đen nhìn cô, định nói gì đó, nhưng rồi khựng lại. Ba chiếc chuông trên cổ tay cô rung lên khẽ khàng, ngân dài giữa không gian chết lặng, như lời cảnh báo mơ hồ về điều gì đó sắp xảy ra.

Cô nàng tóc đỏ lại cúi xuống, cầm cọ lên, giọng cô trôi đi giữa hơi thở đều đặn: ``Có lẽ lần này\dots\ mình sắp vẽ được rồi.''

\end{document}